% Part: incompleteness
% Chapter: representability-in-q
% Section: sigma1-completeness

\documentclass[../../../include/open-logic-section]{subfiles}

\begin{document}

\olfileid[fr]{inc}{inp}{s1c}
\olsection{Complétude \texorpdfstring{$\Sigma_1$}{Sigma-1}}

Malgré l'incomplétude de~$\Th{Q}$ et de ses extensions
cohérentes et axiomatisables, nous avons vu que $\Th{Q}$
démontre de nombreux faits élémentaires sur les numéraux.
Cette capacité va en réalité bien plus loin. Pour en préciser
la portée, introduisons les !!{formula}s $\Delta_0$,
$\Sigma_1$ et $\Pi_1$. Schématiquement, une !!{formula}
$\Sigma_1$ a la forme $\lexists[x][!B(x)]$, où $!B$ se
construit à l'aide des connecteurs propositionnels et de
quantificateurs bornés seulement. Nous montrerons que, si un
!!{sentence}~$!A$ de classe $\Sigma_1$ est vrai dans
$\Struct{N}$, alors $\Th{Q}\Proves !A$
(\olref{thm:sigma1-completeness}).

\begin{defn}
\ollabel{defn:bd-quant}
Une \emph{!!{formula} existentielle bornée} a la forme
$\lexists[x][(x<t\land !A(x))]$, où $t$ est un terme ne
contenant pas librement~$x$. Nous l'écrivons par convention
$\bexists{x<t}{!A(x)}$.
%
Une \emph{!!{formula} universelle bornée} a la forme
$\lforall[x][(x<t\lif !A(x))]$, sous la même condition
sur~$t$. Nous l'écrivons par convention
$\bforall{x<t}{!A(x)}$.
\end{defn}

\begin{defn}
\ollabel{defn:delta0-sigma1-pi1-frm}
Une !!{formula}~$!B$ est $\Delta_0$ si elle se construit à
partir de !!{formula}s atomiques au moyen des seuls
connecteurs propositionnels et quantificateurs bornés.
%
Une !!{formula}~$!A$ est $\Sigma_1$ si
$!A\ident\lexists[x][!B(x)]$ avec $!B$ de classe $\Delta_0$.
%
Une !!{formula}~$!A$ est $\Pi_1$ si
$!A\ident\lforall[x][!B(x)]$ avec $!B$ de classe $\Delta_0$.
\end{defn}

\begin{lem}
\ollabel{lem:q-proves-clterm-id}
Si $t$ est un terme clos et $\Value{t}{N}=n$, alors
$\Th{Q}\Proves\eq[t][\num n]$.
\end{lem}

\begin{proof}
Raisonnons par récurrence sur la complexité de~$t$.
Pour le cas initial, $\Value{\Obj 0}{N}=0$ et
$\Th{Q}\Proves\eq[\Obj 0][\num 0]$, puisque
$\num 0\ident\Obj 0$.
%
Pour l'étape de récurrence, soient $t_1,t_2$ des termes
dont les valeurs sont respectivement $n_1,n_2$, et
supposons
$\Th{Q}\Proves\eq[t_1][\num{n_1}]$ et
$\Th{Q}\Proves\eq[t_2][\num{n_2}]$.

La valeur de $t_1'$ est $n_1+1$. En appliquant les règles
de l'égalité à l'hypothèse de récurrence,
$\Th{Q}\Proves\eq[t_1'][{\num{n_1}}']$.
Par définition des numéraux,
$\num{n_1}'=\num{n_1+1}$; le cas du successeur est établi.

Pour l'addition,
\[
      \Value{(t_1 + t_2)}{N}
    = \Value{t_1}{N} + \Value{t_2}{N}
    = n_1 + n_2.
\]
Les hypothèses de récurrence et la substitution des égaux
donnent d'abord
$\Th{Q}\Proves\eq[t_1+t_2][\num{n_1}+t_2]$,
puis
$\Th{Q}\Proves\eq[t_1+t_2][\num{n_1}+\num{n_2}]$.
Or \olref[inc][req][bre]{lem:q-proves-add} donne
$\Th{Q}\Proves
\eq[\num{n_1}+\num{n_2}][\num{n_1+n_2}]$.
La transitivité de l'égalité conclut.

Le même raisonnement s'applique à la multiplication, à
l'aide du \olref[inc][req][bre]{lem:q-proves-mult}.
%
Ces cas épuisent les termes clos de l'arithmétique :
pour chacun d'eux, si $\Value{t}{N}=n$, alors
$\Th{Q}\Proves\eq[t][\num n]$.
\end{proof}

\begin{prob}
Détaillez le cas de la multiplication dans la preuve du
\olref{lem:q-proves-clterm-id}.
\end{prob}

\begin{lem}
\ollabel{lem:atomic-completeness}
Soient $t_1,t_2$ des termes clos. Alors
\begin{enumerate}
\item Si $\Value{t_1}{N}=\Value{t_2}{N}$,
      $\Th{Q}\Proves\eq[t_1][t_2]$.
\item Si $\Value{t_1}{N}\neq\Value{t_2}{N}$,
      $\Th{Q}\Proves\eq/[t_1][t_2]$.
\item Si $\Value{t_1}{N}<\Value{t_2}{N}$,
      $\Th{Q}\Proves t_1<t_2$.
\item Si $\Value{t_2}{N}\leq\Value{t_1}{N}$,
      $\Th{Q}\Proves\lnot(t_1<t_2)$.
\end{enumerate}
\end{lem}

\begin{proof}
Posons $n=\Value{t_1}{N}$ et $m=\Value{t_2}{N}$.
Le \olref{lem:q-proves-clterm-id} donne
$\Th{Q}\Proves\eq[t_1][\num n]$ et
$\Th{Q}\Proves\eq[t_2][\num m]$.
Si $n=m$, la réflexivité et la transitivité de l'égalité
fournissent $\Th{Q}\Proves\eq[t_1][t_2]$.
Si $n\neq m$, \olref[inc][req][bre]{lem:q-proves-neq}
donne $\Th{Q}\Proves\eq/[\num n][\num m]$;
la substitution des égaux fournit
$\Th{Q}\Proves\eq/[t_1][t_2]$.

Pour les deux cas concernant~$<$, utilisons l'axiome~$!Q_8$ :
pour tous $x,y$,
$x<y\liff\lexists[z][\eq[z'+x][y]]$.
Supposons d'abord $n<m$ et choisissons
$k=m-n-1$. Le terme clos ${\num k}'+\num n$ a la
valeur~$m$; par le premier cas du présent lemme,
$\Th{Q}\Proves\eq[{\num k}'+\num n][\num m]$.
En remplaçant les numéraux par $t_1,t_2$, nous obtenons
$\Th{Q}\Proves\eq[{\num k}'+t_1][t_2]$.
L'introduction de l'existentiel et le sens droit-gauche
de~$!Q_8$ donnent $\Th{Q}\Proves t_1<t_2$.

Supposons maintenant $m\leq n$. Si $m=0$,
\olref[inc][req][min]{lem:less-zero} donne
$\Th{Q}\Proves\lnot(t_1<\num 0)$, et
$\eq[t_2][\num 0]$ permet de remplacer la borne par~$t_2$.
Si $m>0$, écrivons $m=r+1$. D'après le
\olref[inc][req][min]{lem:less-nsucc}, dans~$\Th{Q}$,
$t_1<t_2$ impliquerait
$\eq[t_1][\Obj 0]\lor\dots\lor\eq[t_1][\num r]$.
Or, pour chaque $j\leq r$, nous avons $j<n$ et
\olref[inc][req][bre]{lem:q-proves-neq} donne
$\Th{Q}\Proves\eq/[\num n][\num j]$.
Comme $\Th{Q}\Proves\eq[t_1][\num n]$, aucun des
membres de la disjonction n'est possible. Ainsi
$\Th{Q}\Proves\lnot(t_1<t_2)$.
\end{proof}

\begin{lem}
\ollabel{lem:bounded-quant-equiv}
Supposons que $!A$ soit une !!{formula}, que $t$ soit un
terme clos et que $k=\Value{t}{N}$. Alors
\begin{enumerate}
\item $\Th{Q}\Proves\bforall{x<t}{!A(x)}$ si et seulement si
      $\Th{Q}\Proves !A(\num 0)\land\dots\land
      !A(\num{k-1})$.
\item $\Th{Q}\Proves\bexists{x<t}{!A(x)}$ si et seulement si
      $\Th{Q}\Proves !A(\num 0)\lor\dots\lor
      !A(\num{k-1})$.
\end{enumerate}
Pour $k=0$, la conjonction vide vaut $\ltrue$ et la
disjonction vide vaut $\lfalse$.
\end{lem}

\begin{proof}
Démontrons le cas du quantificateur universel borné.
Si $\Value{t}{N}=0$, le membre gauche est démontrable
dans~$\Th{Q}$ : d'après le
\olref[inc][req][min]{lem:less-zero}, aucun~$x$
n'est inférieur à~$\num 0$. Le membre droit est alors
la conjonction vide $\ltrue$, elle aussi démontrable.

Supposons donc $\Value{t}{N}=k+1$. D'après le
\olref{lem:q-proves-clterm-id}, nous pouvons remplacer
la borne~$t$ par~$\num{k+1}$ dans~$\Th{Q}$.
Si $\Th{Q}\Proves\bforall{x<\num{k+1}}{!A(x)}$,
alors, pour chaque $n\leq k$,
\olref{lem:atomic-completeness} donne
$\Th{Q}\Proves\num n<\num{k+1}$.
La logique fournit donc
$\Th{Q}\Proves !A(\num n)$; en réunissant les
$k+1$ instances, nous obtenons
$\Th{Q}\Proves !A(\num 0)\land\dots\land !A(\num k)$.

Inversement, supposons cette conjonction démontrable
dans~$\Th{Q}$ et raisonnons sous l'hypothèse
$x<\num{k+1}$. D'après le
\olref[inc][req][min]{lem:less-nsucc},
$\eq[x][\num 0]\lor\dots\lor\eq[x][\num k]$.
La substitution des égaux et la disjonction des
instances démontrées donnent~$!A(x)$.
Nous obtenons ainsi
$\bforall{x<\num{k+1}}{!A(x)}$.

Le cas du quantificateur existentiel borné se démontre
de façon analogue.
\end{proof}

\begin{prob}
Détaillez le cas existentiel dans la preuve du
\olref{lem:bounded-quant-equiv}.
\end{prob}

\begin{lem}
\ollabel{lem:delta0-completeness}
Si $!A$ est un !!{sentence}~$\Delta_0$ vrai dans
$\Struct{N}$, alors $\Th{Q}\Proves !A$.
\end{lem}

\begin{proof}
Nous prouvons simultanément que $\Th{Q}$ démontre tout énoncé
$\Delta_0$ vrai dans~$\Struct{N}$ et qu'elle démontre la négation
de tout énoncé $\Delta_0$ faux dans~$\Struct{N}$. La récurrence
porte sur le nombre de connecteurs et de quantificateurs de la
formule. Pour une sous-formule qui comporte des variables libres,
on applique l'hypothèse de récurrence à ses instances closes
obtenues en substituant des numéraux. Cette substitution
ne change pas le nombre de connecteurs et de quantificateurs.

Le cas atomique découle du \olref{lem:atomic-completeness}.
La logique pure donne $\ltrue$ et $\lnot\lfalse$.
Pour les connecteurs propositionnels, on utilise les deux clauses
de l'hypothèse de récurrence. Ainsi, pour une conjonction vraie,
on démontre les deux membres. Pour une conjonction fausse, on
démontre la négation d'un membre faux. Pour une disjonction vraie,
on démontre l'un des membres vrais. Pour une disjonction fausse,
on démontre les négations des deux membres. Pour $\lnot !A$,
on applique l'hypothèse de récurrence à~$!A$. Si $!A$ est vraie,
la preuve de~$!A$ donne celle de~$\lnot\lnot !A$.
Pour l'implication et le biconditionnel, les valeurs de vérité
des deux membres déterminent de même laquelle de la formule
ou de sa négation découle de leurs instances signées.

Pour $\bforall{x<t}{!A(x)}$, la borne~$t$ est close et
$k=\Value{t}{N}$. Si l'énoncé est vrai, chaque instance
$!A(\num n)$ pour $n<k$ est vraie. L'hypothèse de récurrence donne
une preuve de chacune dans~$\Th{Q}$. Le
\olref{lem:bounded-quant-equiv}
donne $\Th{Q}\Proves\bforall{x<t}{!A(x)}$.
Si l'énoncé est faux, une instance est fausse pour un $n<k$.
On démontre sa négation et l'inégalité $\num n<t$, puis la
négation de l'universelle bornée. Le cas de l'existentielle
bornée se traite de même, à l'aide de la disjonction finie.
Lorsque $k=0$, on applique les conventions du
\olref{lem:bounded-quant-equiv} pour la conjonction et la
disjonction vides.
\end{proof}

\begin{prob}
Détaillez le cas existentiel de la preuve du
\olref{lem:delta0-completeness}.
\end{prob}

\begin{thm}
\ollabel{thm:sigma1-completeness}
Si un !!{sentence}~$!A$ de classe $\Sigma_1$ est vrai
dans~$\Struct{N}$, alors $\Th{Q}\Proves !A$.
\end{thm}

\begin{proof}
Écrivons $!A\ident\lexists[x][!B(x)]$, avec~$!B$
de classe $\Delta_0$. Si cet !!{sentence} est vrai
dans~$\Struct{N}$, il existe un entier naturel~$n$
et une assignation~$s$ telle que $s(x)=n$ et
$\Sat{N}{!B(x)}[s]$. Les règles de satisfaction
pour la substitution donnent
$\Sat{N}{!B(\num n)}$.
Cette instance est un !!{sentence}~$\Delta_0$;
par le \olref{lem:delta0-completeness},
$\Th{Q}\Proves !B(\num n)$. L'introduction
de l'existentiel donne enfin
$\Th{Q}\Proves\lexists[x][!B(x)]$, c'est-à-dire
$\Th{Q}\Proves !A$.
\end{proof}

\begin{explain}
Pour un énoncé $\Delta_0$ clos, chaque borne se ramène à un
numéral une fois ses paramètres remplacés. Les quantificateurs
bornés se traitent donc par des preuves portant sur un nombre
fini d'instances. Pour un énoncé $\Sigma_1$ vrai, un seul témoin
standard suffit à produire une telle preuve.
\end{explain}

\end{document}
